\documentclass[12pt,english,a4paper]{article}
\usepackage{amsmath, eurosym, textcomp, graphicx, amsthm, amssymb,amsfonts, natbib, url, xcolor, tikz, subfig,sgame,hyperref,float}

\usepackage[english]{babel} %English language if Koma script classes are used (Bibliography instead of Inhaltsverzeichnis..)
\usepackage[latin1]{inputenc}
\usepackage[T1]{fontenc} %Encoding 
\usepackage{ae,aecompl}%avoids the problem of faint or hardly printable output
\usepackage{pdfsync} %in supporting pdf viewers, this enables forward and inverse search
\synctex=1
%\usepackage{lineno}
%\usepackage{soul} % defines the command \ul{} to underline. Together with the next line the line is set on a level directly under "a" which means that f,g,y will be crossed in their lower parts
%\setuldepth{a}
%\usepackage{setspace} %for varied line spacing
%\onehalfspacing % LaTeX 1.5 line spacing
%\setstretch{1.8} %1.5 line spacing on the Microsoft Word scale
\setlength {\parindent}{0cm}% no indenting for new paragraphs

%\setlength{\parskip}{1.5ex plus 0.5ex minus 0.5ex}%flexible paragraph skipping

\newtheorem{prop}{Proposition}
\newtheorem{theorem}{Theorem}
\newtheorem{assumption}{Assumption}
\newtheorem{lemma}{Lemma}
\newtheorem{proposition}{Proposition}
\newtheorem{observation}{Observation}
\newtheorem{corollary}{Corollary}
\theoremstyle{plain}

\newcommand{\ie}{{\it i.e. }}
\newcommand{\eg}{{\it e.g. }}
\newcommand{\e}{\varepsilon}

\makeatletter
\renewcommand{\section}{\@startsection{section}{1}{0mm}{-1.5\baselineskip}{0.8\baselineskip}{\normalfont\large\centering}}
\renewcommand{\subsection}{\@startsection{subsection}{2}{0mm}{-0.1\baselineskip}{0.5\baselineskip}{\normalfont\bf\flushleft}}
\renewcommand{\@seccntformat}[1]{\csname the#1\endcsname
\hspace{+0mm}\large{.}\hspace{+1.9mm}}
\renewcommand{\@seccntformat}[2]{\csname the#1\endcsname
\hspace{+0mm}\large{.}\hspace{+1.9mm}}
\makeatother

\usepackage[top=2.7cm, bottom=2.7cm, left=2.7cm, right=2.7cm]{geometry}

\def\Re{{\mathbb{R}}}
\def\One{{\mathbb{1}}}

\renewcommand{\baselinestretch}{1.5} %one half spacing

\newenvironment{christoph}{\color{blue}{ }{}} %pretty cool for commenting in color so that co-author can see changes; use \christoph{balbla}

\newenvironment{jan}{\color{red}{ }{}}

%\author{Christoph Schottm\"uller}
\title{Game Theory: Midterm Assignment}
\date{}

\begin{document}
\maketitle

%\setpagewiselinenumbers % line numbers start from 1 on every page
%\modulolinenumbers[5] % only print the line number for every 5th line
%\linenumbers % activate line numbering
\vspace*{-2cm}
\begin{large}
  \begin{center}
    Hand in via email (christoph.schottmueller@econ.ku.dk). \\
    Don't forget to write the names of all group members on the assignment!\\
    Groups consist of up to 3 students.\\
    Deadline: October 30, 2016; 18:00
  \end{center}
\end{large}

\textbf{Exercise 1: }\emph{(decision making under uncertainty)} Explain why the expected utility theorem by von Neumann and Morgenstern is important for game theory. (maximum length: half a page)

% \textbf{Exercise 1: }\emph{(evolutionary stable strategies)} There are two types of creatures: ``dominant'' (d) and ``shy'' (s). We denote the share of dominant types as $\alpha$. The number of expected offsprings of a d type is $0.5*\alpha+1.5*(1-\alpha)$. The number of expected offsprings of a s type is $\alpha*1.5+(1-\alpha)*1$.\\
% Give the corresponding normal form game and determine all Nash equilibria in this game (pure and mixed). Which strategies are evolutionary stable?
%%mixed eq is 1/3 d and 2/3 s; this is stable; two pure eq are not stable

% \textbf{Exercise 2: }\emph{(purification) }Find a game that has a Nash equilibrium in completely mixed strategies (payoffs must not coincide with any game used in the lecture or the exercises, i.e. find your own game!). Using Harsanyi's purification approach, perturb the payoffs in this game and derive a Bayesian Nash equilibrium of the perturbed game. Show that the Bayesian Nash equilibrium approaches the mixed strategy equilibrium when the perturbation goes to zero.

\textbf{Exercise 2: }\emph{(rationalizability) }The luggage of two travelers (A and B) got lost. The baggage of each traveler consisted of one souvenir (and nothing else). The lost souvenirs of A and B are completely identical and have therefore the same value. The airline that lost the luggage has to reimburse the loss. The responsible airline manager does not know the value of the souvenir. He only knows that the value is between 100\$ and 200\$. The travelers know the exact value and they also know that the airline manager only knows that the value is between 100\$ and 200\$. In order to not reimburse too much, the airline manager designs the following game:\\
Each traveler is in a separate room and has to announce the value of the souvenir. If both announce the same value, the airline reimburses the announced value. If the two announcements differ, the traveler with the lower announcement gets his announcement plus 5\$. The traveler with the higher anouncement gets the lower announcement minus 5\$. Announcements must be natural numbers between 100\$ and 200\$ and players care only about the money they expect to get.\\
Is the announcement 200\$ rationalizable?\\
Which announcements are rationalizable?\\
What is/are the Nash equilibrium/equilibria of this game?

\textbf{Exercise 3: }\emph{(correlated equilibrium) } Find a correlated equilibrium in the game below that does not correspond to a Nash equilibrium.
\begin{table}[h]
  \centering
  \begin{game}{3}{3}
    &L&C&R\\
T&0,0&1,2&2,1\\
M&2,1&0,0&1,2\\
B&1,2&2,1&0,0
  \end{game}
\caption{Game for exercise 3}
\end{table}



\textbf{Exercise 4: }\emph{(supermodular games)} $n$ people participate in a contest. The winner of the contest receives a prize of 1\$. All contestants can exert effort $x_i\in[0,1]$. The probability of winning for player $i$ is $\frac{x_i}{\sum_{j=1}^n x_j}$. Each player maximizes his expected prize sum (in \$) minus his effort. Is this game supermodular? 
% No!

% \textbf{Exercise 5: }\emph{(supermodular games)} We look at a pricing game of two firms with zero costs. Each firm sets its price $p_i\in[0,p^m]$ (where $p^m$ is a maximum price). Demand for firm $i$ is $D_i(p_i,p_j)$ which is (i) strictly decreasing in $p_i$, (ii) strictly increasing in $p_j$ and (iii) twice continuously differentiable. Furthermore, the price elasticity of demand (with respect to $p_i$) is a decreasing function of $p_j$. Show that this pricing game is supermodular if we use trick 2 (i.e. letting firms maximize log profits). 

\textbf{Exercise 5: }\emph{(supermodular games)} We use a ``simple search market'' as in the lecture. That is, each of $N$ traders has to decide how much effort $x_i\in [0,1]$ he exerts. If he finds another trader, he gets a payoff of 1. The probability of finding another trader is (different from the lecture!) $\alpha x_i^2\left(\sum_{j\neq i}x_j\right)+\beta x_i$ where $0<\beta< 1$ and $0<\alpha<1$. The costs of effort are $x_i^2$.
\begin{itemize}
\item[(a)]  Show that the game is supermodular.
\item[(b)] Take $N=2$, $\alpha=3/4$ and $\beta=1/2$. Using the results we derived in the lecture, calculate the lowest and highest serially undominated startegy (i.e. $\underline u_i$ and $\bar u_i$). What is the set of rationalizable actions?
\end{itemize}
% foc: $2\alpha x_i x_j+\beta-2 x_i=0$. Note that the payoff function is concave! Using the foc, there are two symmetric pure strategy NE defined by $\alpha x^2+\beta/2-x=0$. The solutions are $x=0.5/\alpha\pm\sqrt{0.25/\alpha^2-\beta/(2\alpha)}$ and therefore $x=1/3$ and $x=1$ for the given parameters. As the game is symmetric and these are the only two symmetric pure strategy NE, we get $\underline u_i=1/3$ and $\bar u_i=1$

\textbf{Exercise 6: }\emph{(existence and perfect equilibrium)} Let $G=\langle N,(A_i),(u_i)\rangle$ be a strategic game in which $N=\{P1,P2\}$, $A_1=\{T,B\}$ and $A_2=\{L,R\}$, i.e. $G$ is a $2\times 2$ game. Similar to the definition of perfect equilibrium, we define a game $G'(\eta)$: We require that in $G'$ each player plays a completely mixed strategy  and plays each of his actions with at least probability $\eta\in(0,1/2)$. Everything else is as in $G$. Use the Brouwer fixed point theorem to show that an equilibrium exists in $G'$.

\textbf{Exercise 7: }\emph{ } Consider the game in the table below. Which outcome do you predict? What would the solution concepts (or refinements) covered in the course predict?

\begin{table}[h!]
  \centering
  \begin{game}{3}{3}
    &L   &C   &R\\
U   &  0,3&0,1&1,2\\
M   &  1,1&2,0&0,1\\
D   &  1,1&1,2&3,3
  \end{game}
\label{tab:1}
\caption{Game for exercise 7}
\end{table}

\textbf{Exercise 8: }\emph{(perfect equilibrium)} Determine the two pure strategy Nash equilibria in the following game. Show that one of the Nash equilibria is perfect and the other is not.
\begin{table}[h!]
  \centering
  \begin{game}{3}{3}
    &L   &C   &R\\
U   &  0,3&0,1&3,3\\
M   &  2,2&2,0&0,1\\
D   &  1,1&1,2&3,0
  \end{game}
\label{tab:1}
\caption{Game for exercise 8}
\end{table}

% \textbf{Exercise 7: }\emph{(knowledge) }Think of the hat game we had in the lecture: Each of $n$ people in a room wear a hat (you can assume that $n$ is even). They can see the colour of everyone else's hat but not the colour of their own one. Hats are either black or white. A person raises his hand if he knows the colour of his hat.\\
% An outside observer (who does not wear a hat) announces the following: \textbf{``At least half of you wear a white hat.''} It is common knowledge that the observer told the truth.\\
% What is common knowledge after this announcement? Assume that no one raises his hand after this announcement. What is common knowledge then? Now the outside observer counts rounds (after his announcement). What is common knowledge after round $k$ if no one raised his hand until then? After how many rounds will someone raise his hand at the latest?

\textbf{Exercise 9: }Read the paper ``Rationalizable conjectural equilibrium: between Nash and rationalizability'' (\url{http://dx.doi.org/10.1006/game.1994.1016} to download it for free conncect from campus or via VPN). Answer the following three questions in (together!) less than 1 page:  What is the main idea of the paper? Why does it make sense or not make sense? How does it relate to concepts you learned in the course?





%\appendix
%\setcounter{table}{0}
%\renewcommand{\thetable}{\thesection\arabic{table}}
%
%\newpage
%
%\begin{center}
%\large{Appendix}
%\end{center}
%
%
%\newpage
%\bibliographystyle{chicago}
%\bibliography{}


\end{document}
